%% Drafted from the mapped corpus by scripts/draft_topics.py.
% Model: gpt-5.6-luna; source characters: 9356.
\begin{konspekt}
\kreq{Официално заглавие}{„Структури от данни. Стек, опашка, списък, дърво. \emph{Основни операции върху тях. Реализация}.“}
\kselect{За изпита ще бъдат избрани \textbf{две} от посочените структури, които да бъдат описани. Възможни са и задачи за \emph{практическата} част на изпита.}
\kreq{Дефиниция}{на понятието структура от данни --- \kref{18.1}.}
\kreq{За всяка структура се иска еднакъв набор}{логическо описание; характеристики на реализациите; сложност на операциите; \textbf{дефиниране на клас}, използващ една от реализациите.}
\kreq{Реализации}{статична, динамична и свързана --- назованите от конспекта три категории --- \kref{18.2}.}
\kreq{Списък}{с една и две връзки; сложност на добавяне, премахване и намиране --- \kr{18.5}; клас --- \kref{18.5.3}.}
\kreq{Стек}{логическо описание и представяния --- \kr{18.3}; клас --- \kref{18.3.3}.}
\kreq{Опашка}{логическо описание и представяния --- \kr{18.4}; клас --- \kref{18.4.3}.}
\kreq{Дървовидни структури}{кореново и двоично кореново дърво; начини за представяне в паметта --- \kr{18.6}; клас --- \kref{18.6.4}.}
\kreq{Двоично кореново дърво за търсене}{логическо описание, представяне, сложност --- \kr{18.7}; клас --- \kref{18.7.5}.}
\kreq{Обобщение}{таблица със сложностите и допусканията --- \kref{18.8.1}.}
\kreq{Литература}{[11], [15], [25], [28].}
\end{konspekt}

\section{Структури от данни}\label{s:18.1}

\begin{defn}[Структура от данни]
Структура от данни е организирана информация, която може да бъде описана, създадена и обработвана с помощта на програма.
\end{defn}

При определяне на една структура от данни се разграничават два аспекта:

\begin{itemize}
    \item логическо описание на структурата -- описание чрез нейната декомпозиция, чрез по-прости структури и чрез операциите над нея, сведени до по-прости операции;
    \item физическо представяне -- методът, по който структурата се разполага в паметта на компютъра.
\end{itemize}

Логическата структура определя какви са елементите, как са свързани и какви операции са допустими. Физическото представяне определя как тази информация се реализира чрез клетки памет, масиви, указатели и други програмни средства.

В настоящата глава се разглеждат четири основни структури от данни: стек, опашка, списък и дърво.

\section{Трите категории реализации}\label{s:18.2}

Конспектът различава \emph{статична}, \emph{динамична} и \emph{свързана}
реализация за стека и за опашката. Разликата е в това как се управлява паметта,
а не в логическото описание, което остава едно и също.

\begin{defn}[Статична (последователна) реализация]
Елементите се пазят в масив с \emph{фиксиран} капацитет $N$, заделен предварително.
Достъпът по индекс е константен, но при $N$ елемента следващото включване е
невъзможно: настъпва \emph{препълване}. Паметта е заета изцяло независимо от
броя на реално съхранените елементи.
\end{defn}

\begin{defn}[Динамична реализация (масив с преоразмеряване)]
Елементите се пазят в масив, заделен в динамичната памет, чийто капацитет се
\emph{удвоява}, когато се запълни: заделя се нов блок с двоен размер, елементите
се копират и старият блок се освобождава. Няма предварителна граница за броя
елементи; достъпът по индекс остава константен.
\end{defn}

\begin{prop}[Амортизирана цена на включването]
При удвояване на капацитета всяко включване струва $O(1)$ \emph{амортизирано},
макар отделно включване да струва $O(n)$, когато предизвика преоразмеряване.
\end{prop}

\begin{proof}
Разглеждаме редица от $n$ включвания в стек, започнал с капацитет $c_0\ge 1$.
Преоразмеряванията стават при размери $c_0,2c_0,4c_0,\ldots$, а цената на всяко е
пропорционална на броя копирани елементи. Общата цена на копиранията е
\[
c_0+2c_0+4c_0+\cdots+2^kc_0
<2\cdot 2^kc_0\le 2n\cdot\frac{2^kc_0}{n},
\]
където $2^kc_0\le n$ е последният капацитет преди спиране; следователно сумата е
по-малка от $2n$. Заедно с $n$-те константни записа общата цена е $O(n)$, тоест
$O(1)$ на включване средно.
\end{proof}

\begin{defn}[Свързана реализация]
Всеки елемент се пази в отделен възел с указател към следващия. Няма
преоразмеряване и няма препълване, но за всеки елемент се плаща допълнителна
памет за указателя, а достъпът до $i$-тия елемент е $O(i)$, защото се минава
през веригата.
\end{defn}

\begin{remark}
Изборът не е въпрос на вкус. Статичната реализация се използва при известна
горна граница и при забрана за динамична памет; динамичната --- когато е нужен
пряк достъп по индекс без такава граница; свързаната --- когато преобладават
включвания и изключвания в известна позиция.
\end{remark}

\section{Стек}\label{s:18.3}

\begin{defn}[Стек]
Стекът е крайна редица от елементи от един и същи тип, при която включването и изключването на елементи са допустими само в единия край на редицата. Този край се нарича връх на стека.
\end{defn}

Стекът е линейна динамична структура от данни. При него е възможен пряк достъп само до елемента на върха. За останалите елементи достъпът се осъществява последователно, започвайки от върха.

\begin{keydefn}[LIFO]
Правилото LIFO (\emph{Last In, First Out}) означава, че последно включеният в стека елемент е първият изключен от него.
\end{keydefn}

Основните операции върху стек са:

\begin{itemize}
    \item \emph{включване} (\emph{push}) -- добавяне на елемент на върха;
    \item \emph{изключване} (\emph{pop}) -- премахване на елемента от върха;
    \item \emph{достъп до върха} -- прочитане на елемента на върха без премахването му;
    \item \emph{проверка за празен стек};
    \item \emph{търсене} на елемент.
\end{itemize}

Включването и изключването на елемент са константни операции, тоест имат сложност $O(1)$. Търсенето на елемент има линейна сложност $O(n)$, тъй като в най-общия случай трябва да се прегледат всички елементи.

\subsection{Последователно представяне на стек}\label{s:18.3.1}

При последователното, или статичното, представяне предварително се заделя блок от памет, обикновено реализиран чрез масив. В този блок стекът нараства и се свива. Поддържа се индекс или указател към върха.

Ако масивът е с размер $N$, стекът може да бъде описан чрез масива \verb|data| и целочислен индекс \verb|top|. При празен стек \verb|top| има стойност, която показва, че няма валиден елемент. При включване стойността на \verb|top| се увеличава и новият елемент се записва на съответната позиция. При изключване елементът на \verb|top| се извлича, след което \verb|top| се намалява.

Примерна реализация на стек от цели числа чрез масив е:

\begin{verbatim}
const int Capacity = 100;

struct Stack {
  int data[Capacity];
  int top;
};

void init(Stack& s) {
  s.top = -1;
}

bool empty(const Stack& s) {
  return s.top == -1;
}

bool full(const Stack& s) {
  return s.top == Capacity - 1;
}

bool push(Stack& s, int value) {
  if (full(s)) {
    return false;
  }

  ++s.top;
  s.data[s.top] = value;
  return true;
}

bool pop(Stack& s, int& value) {
  if (empty(s)) {
    return false;
  }

  value = s.data[s.top];
  --s.top;
  return true;
}
\end{verbatim}

При тази реализация възниква състояние \emph{препълване}, ако се направи опит за включване в стек, който е достигнал капацитета си. При изключване от празен стек възниква състояние \emph{празен стек} или неуспешна операция. Размерът на последователно представения стек е ограничен от предварително заделения блок.

\subsection{Свързано представяне на стек}\label{s:18.3.2}

При свързаното представяне последователните елементи на стека могат да се намират на различни места в паметта. Връзката между два съседни елемента се осъществява чрез указател към следващия елемент. Указателят на последния елемент обикновено има стойност \verb|nullptr|.

За задаване на стека е достатъчен указател към неговия връх. Един елемент на стек от цели числа може да има вида:

\begin{verbatim}
struct StackItem {
  int data;
  StackItem* link;
};
\end{verbatim}

Полето \verb|data| съдържа информацията, а полето \verb|link| съдържа адреса на следващия елемент към вътрешността на стека.

\begin{verbatim}
struct LinkedStack {
  StackItem* top;
};

void init(LinkedStack& s) {
  s.top = nullptr;
}

bool empty(const LinkedStack& s) {
  return s.top == nullptr;
}

void push(LinkedStack& s, int value) {
  StackItem* item = new StackItem;
  item->data = value;
  item->link = s.top;
  s.top = item;
}

bool pop(LinkedStack& s, int& value) {
  if (empty(s)) {
    return false;
  }

  StackItem* item = s.top;
  value = item->data;
  s.top = item->link;
  delete item;
  return true;
}
\end{verbatim}

При включване новият елемент се свързва с досегашния връх и става нов връх. При изключване се запазва стойността на върха, указателят се премества към следващия елемент, а старият елемент се освобождава. И двете операции са с константна сложност.

Свързаното представяне не изисква предварително зададен максимален брой елементи, но за всеки елемент се съхранява и указател. То е подходящо, когато размерът на стека се изменя динамично.


\subsection{Клас за стек (динамична реализация)}\label{s:18.3.3}

Конспектът иска \emph{дефиниране на клас} за стек, използващ една от
реализациите. Функциите по-горе работят върху отворена структура: всеки може да
пипне \verb|top| и да наруши съгласуваността. Класът затваря представянето и
поема отговорност за паметта.

Инвариант на класа: \verb|data| сочи блок от точно \verb|capacity| клетки в
динамичната памет, $0\le$ \verb|size| $\le$ \verb|capacity|, валидни са
елементите \verb|data[0..size-1]|, а върхът е \verb|data[size-1]|.

\begin{verbatim}
#include <stdexcept>

class Stack {
  int* data;       // блок с capacity клетки
  int size;        // брой съхранени елементи
  int capacity;    // размер на блока

  void resize(int newCapacity) {
    int* moved = new int[newCapacity];
    for (int i = 0; i < size; ++i) moved[i] = data[i];
    delete[] data;
    data = moved;
    capacity = newCapacity;
  }

public:
  Stack() : data(new int[4]), size(0), capacity(4) {}

  Stack(const Stack& other)
      : data(new int[other.capacity]),
        size(other.size), capacity(other.capacity) {
    for (int i = 0; i < size; ++i) data[i] = other.data[i];
  }

  Stack& operator=(const Stack& other) {
    if (this != &other) {
      int* moved = new int[other.capacity];
      for (int i = 0; i < other.size; ++i)
        moved[i] = other.data[i];
      delete[] data;
      data = moved;
      size = other.size;
      capacity = other.capacity;
    }
    return *this;
  }

  ~Stack() { delete[] data; }

  bool empty() const { return size == 0; }

  void push(int value) {
    if (size == capacity) resize(2 * capacity);
    data[size++] = value;
  }

  int pop() {
    if (empty())
      throw std::underflow_error("pop от празен стек");
    return data[--size];
  }

  int top() const {
    if (empty())
      throw std::underflow_error("top от празен стек");
    return data[size - 1];
  }
};
\end{verbatim}

Трите члена --- копиращ конструктор, операция за присвояване и деструктор ---
вървят заедно. Ако липсват, копирането на \verb|Stack| копира указателя, двата
обекта освобождават един и същи блок и програмата се разрушава. Това е
управлението на ресурси, което Въпрос 16 нарича RAII: заделянето става в
конструктора, освобождаването --- в деструктора, и никой път през кода не може
да ги пропусне.

Същият клас в \emph{статична} реализация се получава така: \verb|data| става
масив с фиксиран капацитет, \verb|resize| отпада, а \verb|push| хвърля
\verb|std::overflow_error| при пълен стек. В \emph{свързана} реализация трите
полета се заменят с единствен указател \verb|top| към възел, както в \S18.3.2.

\section{Опашка}\label{s:18.4}

\begin{defn}[Опашка]
Опашката е крайна редица от елементи от един и същи тип, при която включването е допустимо в единия край, а изключването -- в другия край.
\end{defn}

Краят, в който се добавят елементи, се нарича \emph{край на опашката}, а краят, от който се премахват елементи, се нарича \emph{начало на опашката}. Пряк достъп е възможен само до елемента в началото на опашката.

\begin{keydefn}[FIFO]
Правилото FIFO (\emph{First In, First Out}) означава, че първият включен в опашката елемент е първият изключен от нея.
\end{keydefn}

Основните операции върху опашка са:

\begin{itemize}
    \item включване на елемент в края;
    \item изключване на елемент от началото;
    \item достъп до елемента в началото;
    \item проверка за празна опашка;
    \item търсене на елемент.
\end{itemize}

Добавянето и премахването на елемент са операции с константна сложност $O(1)$. Търсенето, когато е разрешено, е линейно по броя на елементите.

\subsection{Последователно представяне на опашка}\label{s:18.4.1}

При последователното представяне предварително се заделя блок памет, реализиран чрез масив. Обикновено блокът се разглежда като цикличен. Това означава, че след достигане на последната позиция следва първата позиция, ако в началото има освободено място.

Необходимо е да се пазят поне два индекса:

\begin{itemize}
    \item индекс на началото на опашката;
    \item индекс на следващата свободна позиция в края на опашката.
\end{itemize}

Необходимо е също да се знае броят на елементите или да се използва допълнителен признак, който различава празна от пълна опашка.

\begin{verbatim}
const int Capacity = 100;

struct Queue {
  int data[Capacity];
  int first;
  int last;
  int count;
};

void init(Queue& q) {
  q.first = 0;
  q.last = 0;
  q.count = 0;
}

bool empty(const Queue& q) {
  return q.count == 0;
}

bool full(const Queue& q) {
  return q.count == Capacity;
}

bool enqueue(Queue& q, int value) {
  if (full(q)) {
    return false;
  }

  q.data[q.last] = value;
  q.last = (q.last + 1) % Capacity;
  ++q.count;
  return true;
}

bool dequeue(Queue& q, int& value) {
  if (empty(q)) {
    return false;
  }

  value = q.data[q.first];
  q.first = (q.first + 1) % Capacity;
  --q.count;
  return true;
}
\end{verbatim}

При включване елементът се записва на позицията, определена от \verb|last|, след което индексът се придвижва циклично. При изключване се извлича елементът на позиция \verb|first| и началният индекс също се придвижва циклично.

\subsection{Свързано представяне на опашка}\label{s:18.4.2}

Свързаното представяне на опашка е подобно на свързаното представяне на стек. Елементите се намират на различни места в паметта и са свързани чрез указатели. За разлика от стека, при опашката са нужни два указателя:

\begin{itemize}
    \item указател към началото на опашката;
    \item указател към края на опашката.
\end{itemize}

Елементът на опашката може да бъде описан така:

\begin{verbatim}
struct QueueItem {
  int data;
  QueueItem* link;
};
\end{verbatim}

При празна опашка и двата указателя имат стойност \verb|nullptr|. При първото включване и двата трябва да бъдат насочени към новия елемент. При следващи включвания новият елемент се свързва след досегашния край, а указателят към края се премества към него. При изключване се премахва елементът от началото и указателят към началото се премества към следващия елемент. Ако след премахването опашката стане празна, указателят към края също се установява на \verb|nullptr|.

\begin{verbatim}
struct LinkedQueue {
  QueueItem* first;
  QueueItem* last;
};

void init(LinkedQueue& q) {
  q.first = nullptr;
  q.last = nullptr;
}

bool empty(const LinkedQueue& q) {
  return q.first == nullptr;
}

void enqueue(LinkedQueue& q, int value) {
  QueueItem* item = new QueueItem;
  item->data = value;
  item->link = nullptr;

  if (q.last == nullptr) {
    q.first = item;
    q.last = item;
  } else {
    q.last->link = item;
    q.last = item;
  }
}

bool dequeue(LinkedQueue& q, int& value) {
  if (empty(q)) {
    return false;
  }

  QueueItem* item = q.first;
  value = item->data;
  q.first = item->link;

  if (q.first == nullptr) {
    q.last = nullptr;
  }

  delete item;
  return true;
}
\end{verbatim}

При наличие на указател към края включването е константна операция. Без такъв указател намирането на последния елемент би изисквало преминаване през опашката.


\subsection{Клас за опашка (свързана реализация)}\label{s:18.4.3}

Инвариант: или \verb|head == nullptr| и \verb|tail == nullptr| (празна опашка),
или \verb|head| сочи първия възел, \verb|tail| --- последния,
\verb|tail->next == nullptr|, и всеки възел е достижим от \verb|head|.

\begin{verbatim}
#include <stdexcept>

class Queue {
  struct Node {
    int data;
    Node* next;
  };

  Node* head;   // оттук се изключва
  Node* tail;   // тук се включва

  void clear() {
    while (head != nullptr) {
      Node* item = head;
      head = head->next;
      delete item;
    }
    tail = nullptr;
  }

  void copyFrom(const Queue& other) {
    head = tail = nullptr;
    for (Node* it = other.head; it != nullptr; it = it->next)
      enqueue(it->data);
  }

public:
  Queue() : head(nullptr), tail(nullptr) {}
  Queue(const Queue& other) { copyFrom(other); }

  Queue& operator=(const Queue& other) {
    if (this != &other) {
      clear();
      copyFrom(other);
    }
    return *this;
  }

  ~Queue() { clear(); }

  bool empty() const { return head == nullptr; }

  void enqueue(int value) {
    Node* item = new Node{value, nullptr};
    if (tail != nullptr) tail->next = item;
    else head = item;
    tail = item;
  }

  int dequeue() {
    if (empty())
      throw std::underflow_error("dequeue от празна опашка");
    Node* item = head;
    int value = item->data;
    head = head->next;
    if (head == nullptr) tail = nullptr;
    delete item;
    return value;
  }

  int front() const {
    if (empty())
      throw std::underflow_error("front от празна опашка");
    return head->data;
  }
};
\end{verbatim}

Двата указателя са това, което прави и включването, и изключването $O(1)$: без
\verb|tail| включването в края би изисквало обхождане на цялата верига.
Присвояването на \verb|tail = nullptr| при изпразване не е козметика --- без
него \verb|tail| остава да сочи освободена памет.

\section{Списък}\label{s:18.5}

\begin{defn}[Списък]
Списъкът е крайна редица от елементи от един и същи тип, при която включването и изключването са допустими от произволно място в редицата.
\end{defn}

За разлика от стека и опашката списъкът позволява работа с произволна позиция. Възможен е достъп до всички елементи, който може да бъде пряк или непряк в зависимост от физическото представяне.

Разглеждат се три основни начина за представяне на списък:

\begin{itemize}
    \item последователно представяне;
    \item едносвързано представяне;
    \item двусвързано представяне.
\end{itemize}

Последователното представяне не се използва като основен начин за списък, тъй като включването и изключването в произволно място обикновено изискват преместване на елементи.

\subsection{Едносвързан списък}\label{s:18.5.1}

При едносвързаното представяне елементите се съхраняват на различни места в паметта. Всеки елемент съдържа информация и указател към следващия елемент. Указателят на последния елемент е \verb|nullptr|.

За задаване на списъка е достатъчен указател към началото. Ако се добави и указател към края, включването в края може да се извършва за константно време.

\begin{verbatim}
struct ListItem {
  int data;
  ListItem* next;
};

struct List {
  ListItem* first;
  ListItem* last;
};
\end{verbatim}

Основните операции са:

\begin{itemize}
    \item включване в началото;
    \item включване в края;
    \item включване след даден елемент;
    \item изключване от началото;
    \item изключване след даден елемент;
    \item търсене;
    \item обхождане на списъка.
\end{itemize}

Включването и изключването в началото са константни операции. Включването в края също е константно, ако се поддържа указател към последния елемент. Включването или изключването в произволна позиция, когато позицията трябва първо да бъде намерена, е линейно. Търсенето също е линейна операция.

\begin{verbatim}
void insertFirst(List& list, int value) {
  ListItem* item = new ListItem;
  item->data = value;
  item->next = list.first;
  list.first = item;

  if (list.last == nullptr) {
    list.last = item;
  }
}

bool removeFirst(List& list, int& value) {
  if (list.first == nullptr) {
    return false;
  }

  ListItem* item = list.first;
  value = item->data;
  list.first = item->next;

  if (list.first == nullptr) {
    list.last = nullptr;
  }

  delete item;
  return true;
}

void insertLast(List& list, int value) {
  ListItem* item = new ListItem;
  item->data = value;
  item->next = nullptr;

  if (list.last == nullptr) {
    list.first = item;
    list.last = item;
  } else {
    list.last->next = item;
    list.last = item;
  }
}
\end{verbatim}

За изключване на елемент, който не е първи, е необходим указател към предходния елемент. След промяната указателят на предходния елемент трябва да прескочи изключвания елемент и да сочи към следващия.

\subsection{Двусвързан списък}\label{s:18.5.2}

При двусвързания списък всеки елемент има два указателя: към следващия и към предходния елемент. Елементът може да бъде описан чрез:

\begin{verbatim}
struct DLListItem {
  int data;
  DLListItem* next;
  DLListItem* prev;
};

struct DLList {
  DLListItem* first;
  DLListItem* last;
};
\end{verbatim}

За списъка обикновено се пазят указатели към началото и края. Допълнително може да се пази указател към текущия елемент, което улеснява някои операции по вмъкване, изключване и търсене.

При едносвързан списък движението по връзките е само напред. При двусвързан списък може да се преминава и в двете посоки. Изключването на текущ елемент се извършва чрез свързване на неговия предходен и следващ елемент един с друг. При крайни елементи единият от съответните указатели е \verb|nullptr|.

\begin{verbatim}
void insertAfter(DLList& list,
                 DLListItem* current,
                 int value) {
  if (current == nullptr) {
    return;
  }

  DLListItem* item = new DLListItem;
  item->data = value;
  item->prev = current;
  item->next = current->next;

  if (current->next != nullptr) {
    current->next->prev = item;
  } else {
    list.last = item;
  }

  current->next = item;
}
\end{verbatim}

Двусвързаното представяне изисква повече памет заради втория указател, но улеснява изключването и движението в обратна посока.


\subsection{Клас за списък (двусвързана реализация)}\label{s:18.5.3}

Инвариант: \verb|count| е броят на възлите; предшественикът на \verb|first| и
наследникът на \verb|last| са \verb|nullptr|; за всеки възел $v$ с наследник е
изпълнено \verb|v->next->prev == v|; при празен списък и двата указателя са
\verb|nullptr|.

\begin{verbatim}
class List {
  struct Node {
    int data;
    Node* prev;
    Node* next;
  };

  Node* first;
  Node* last;
  int count;

  Node* nodeAt(int pos) const {      // изисква 0 <= pos < count
    Node* it = first;
    for (int i = 0; i < pos; ++i) it = it->next;
    return it;
  }

  void clear() {
    while (first != nullptr) {
      Node* item = first;
      first = first->next;
      delete item;
    }
    last = nullptr;
    count = 0;
  }

  void copyFrom(const List& other) {
    first = last = nullptr;
    count = 0;
    for (Node* it = other.first; it != nullptr; it = it->next)
      pushBack(it->data);
  }

public:
  List() : first(nullptr), last(nullptr), count(0) {}
  List(const List& other) { copyFrom(other); }

  List& operator=(const List& other) {
    if (this != &other) {
      clear();
      copyFrom(other);
    }
    return *this;
  }

  ~List() { clear(); }

  int size() const { return count; }
  bool empty() const { return count == 0; }

  void pushBack(int value) {
    Node* item = new Node{value, last, nullptr};
    if (last != nullptr) last->next = item;
    else first = item;
    last = item;
    ++count;
  }

  void insertAt(int pos, int value) {  // изисква 0 <= pos <= count
    if (pos == count) {
      pushBack(value);
      return;
    }
    Node* next = nodeAt(pos);
    Node* item = new Node{value, next->prev, next};
    if (next->prev != nullptr) next->prev->next = item;
    else first = item;
    next->prev = item;
    ++count;
  }

  void removeAt(int pos) {           // изисква 0 <= pos < count
    Node* item = nodeAt(pos);
    if (item->prev != nullptr) item->prev->next = item->next;
    else first = item->next;
    if (item->next != nullptr) item->next->prev = item->prev;
    else last = item->prev;
    delete item;
    --count;
  }

  // индекс на първото срещане или -1
  int find(int value) const {
    int i = 0;
    for (Node* it = first; it != nullptr; it = it->next, ++i)
      if (it->data == value) return i;
    return -1;
  }
};
\end{verbatim}

Двете връзки правят премахването на \emph{известен} възел константно: при
едносвързан списък трябва да се търси предшественикът, тоест $O(n)$. Цената на
\verb|insertAt| и \verb|removeAt| тук е $O(\mathit{pos})$ --- разходът е в
намирането на позицията, не в пренавързването, което е $O(1)$.

\section{Дърво}\label{s:18.6}

\begin{defn}[Дърво]
Дърво е свързан ацикличен граф.
\end{defn}

Дървото е разклонена структура от данни. Елементите му се наричат \emph{върхове} или \emph{възли}. Един от върховете се разглежда като корен. Всеки връх може да има произволен брой наследници. Връзките между върховете определят йерархична структура.

Рекурсивното определение на дърво е:

\begin{itemize}
    \item празното дърво е дърво;
    \item дървото с единствен връх, наречен корен, е дърво;
    \item дърво с корен и произволен брой наследници, всеки от които е корен на дърво, също е дърво.
\end{itemize}

Предоставя се пряк достъп до корена, а до останалите върхове достъпът се осъществява чрез връзките в структурата.

\subsection{Двойно дърво}\label{s:18.6.1}

\begin{defn}[Двойно дърво]
Двойното дърво от тип $T$ е рекурсивна структура от данни, която е или празна, или е образувана от данна от тип $T$, наречена корен, двойно дърво от тип $T$, наречено ляво поддърво, и двойно дърво от тип $T$, наречено дясно поддърво.
\end{defn}

Множеството от върховете на двойното дърво се определя рекурсивно. Празното двойно дърво няма върхове. Непразното двойно дърво има корен и всички върхове на лявото и дясното поддърво.

Левият наследник на връх е коренът на лявото му поддърво, ако такова поддърво съществува. Аналогично се определя десният наследник.

\begin{defn}[Листа и вътрешни върхове]
Върхове без наследници се наричат \emph{листа}. Вътрешни върхове са върховете с поне едно дете, включително корена, когато той има дете.
\end{defn}

Свързаното представяне на двойно дърво се реализира чрез указател към структура с три полета: информационно поле и два адреса за лявото и дясното поддърво.

\begin{verbatim}
struct Node {
  int data;
  Node* left;
  Node* right;
};
\end{verbatim}

Празното поддърво се представя с указател \verb|nullptr|.

\subsection{Операции върху двойно дърво}\label{s:18.6.2}

Основните операции върху двойно дърво са:

\begin{itemize}
    \item достъп до връх;
    \item включване на връх;
    \item изключване на връх;
    \item обхождане.
\end{itemize}

До корена се осъществява пряк достъп, а до останалите върхове -- непряк чрез лявата и дясната връзка. Включването и изключването могат да се извършват на произволно място, като резултатът трябва да бъде двойно дърво от същия тип.

\begin{defn}[Обхождане]
Обхождането е метод, чрез който се осъществява достъп до всеки връх на дървото точно един път.
\end{defn}

При рекурсивното обхождане се изпълняват три действия в определен ред:

\begin{enumerate}
    \item обхождане на корена;
    \item обхождане на лявото поддърво;
    \item обхождане на дясното поддърво.
\end{enumerate}

В зависимост от реда на тези действия се получават различни обходи. Когато редът е корен, ляво поддърво, дясно поддърво, се получава предварително обхождане. Когато първо се обхожда лявото поддърво, след това коренът, а накрая дясното поддърво, се получава смесено обхождане. При обхождане ляво поддърво, дясно поддърво, корен се получава последващо обхождане.

\begin{verbatim}
void preorder(Node* tree) {
  if (tree == nullptr) {
    return;
  }

  visit(tree->data);
  preorder(tree->left);
  preorder(tree->right);
}

void inorder(Node* tree) {
  if (tree == nullptr) {
    return;
  }

  inorder(tree->left);
  visit(tree->data);
  inorder(tree->right);
}

void postorder(Node* tree) {
  if (tree == nullptr) {
    return;
  }

  postorder(tree->left);
  postorder(tree->right);
  visit(tree->data);
}
\end{verbatim}

\subsection{Представяния на двойно дърво}\label{s:18.6.3}

Използват се три начина за представяне на двойно дърво:

\begin{enumerate}
    \item свързано представяне;
    \item верижно представяне чрез масиви;
    \item представяне чрез списък на бащите.
\end{enumerate}

При свързаното представяне всеки възел съдържа стойност и два указателя. Този начин е естествен за рекурсивните операции.

При верижното представяне се използват три масива $a[N]$, $b[N]$ и $c[N]$, където $N$ е броят на върховете. Върховете са номерирани от $0$ до $N-1$. В $a[i]$ се съхранява стойността на $i$-тия връх, в $b[i]$ -- индексът на левия наследник, а в $c[i]$ -- индексът на десния наследник. Ако съответният наследник липсва, се записва $-1$. Отделно се пази индексът на корена.

При представяне чрез списък на бащите се използва един масив $p[N]$. В $p[i]$ се съхранява индексът на единствения баща на върха $i$. За корена се записва $-1$.

\begin{supp}[Бележка]
За общо двойно дърво източникът посочва сложност $O(n)$ за търсене, а на друго място посочва $O(\log n)$ за добавяне. Тези оценки зависят от конкретната структура и от начина на търсене. При обикновено двойно дърво без свойство за подреждане търсенето в най-общия случай е линейно. Сложността $O(\log n)$ е характерна за операции в дърво с подходяща височина, например балансирано двойно подредено дърво.
\end{supp}


\subsection{Клас за двойно кореново дърво}\label{s:18.6.4}

Класът притежава корена и отговаря за освобождаването на цялото дърво.
Инвариант: \verb|root| е \verb|nullptr| (празно дърво) или сочи възел, чиито
поддървета са непресичащи се --- никой възел не е достижим по два различни пътя,
тоест структурата е дърво, а не общ граф.

\begin{verbatim}
#include <iostream>

class BinaryTree {
  struct Node {
    int data;
    Node* left;
    Node* right;
  };

  Node* root;

  static void destroy(Node* v) {
    if (v == nullptr) return;
    destroy(v->left);
    destroy(v->right);
    delete v;
  }

  static Node* clone(const Node* v) {
    if (v == nullptr) return nullptr;
    return new Node{v->data, clone(v->left), clone(v->right)};
  }

  static int heightOf(const Node* v) {
    if (v == nullptr) return 0;
    int l = heightOf(v->left);
    int r = heightOf(v->right);
    return 1 + (l > r ? l : r);
  }

  static int countOf(const Node* v) {
    if (v == nullptr) return 0;
    return 1 + countOf(v->left) + countOf(v->right);
  }

  static void inorder(const Node* v) {
    if (v == nullptr) return;
    inorder(v->left);
    std::cout << v->data << ' ';
    inorder(v->right);
  }

public:
  BinaryTree() : root(nullptr) {}

  // Рекурсивното определение, изразено като конструктор.
  BinaryTree(int value,
             const BinaryTree& left, const BinaryTree& right)
      : root(new Node{value,
                      clone(left.root), clone(right.root)}) {}

  BinaryTree(const BinaryTree& other)
      : root(clone(other.root)) {}

  BinaryTree& operator=(const BinaryTree& other) {
    if (this != &other) {
      Node* copied = clone(other.root);   // първо копираме,
      destroy(root);                      // после освобождаваме
      root = copied;
    }
    return *this;
  }

  ~BinaryTree() { destroy(root); }

  bool empty() const { return root == nullptr; }
  int size() const { return countOf(root); }
  int height() const { return heightOf(root); }
  void printInorder() const { inorder(root); }
};
\end{verbatim}

Конструкторът с три аргумента е точно рекурсивната дефиниция: стойност, ляво
поддърво, дясно поддърво. Затова
\begin{verbatim}
BinaryTree leaf4(4, BinaryTree(), BinaryTree());
BinaryTree leaf12(12, BinaryTree(), BinaryTree());
BinaryTree t(10, leaf4, leaf12);
\end{verbatim}
строи дървото с корен $10$ и наследници $4$ и $12$, без нито един \verb|new| в
кода на потребителя.

В \verb|operator=| копирането предхожда освобождаването: ако \verb|clone| хвърли
изключение поради недостиг на памет, обектът остава непроменен вместо да остане
с освободен корен.

\section{Двойно подредено дърво}\label{s:18.7}

\begin{defn}[Двойно подредено дърво]
Двойно подредено дърво от тип $T$ е двойно дърво, за което върховете на лявото поддърво са по-малки от корена, върховете на дясното поддърво са по-големи от корена, а лявото и дясното поддърво също са двойно подредени дървета от тип $T$.
\end{defn}

Празното двойно дърво се приема за двойно подредено дърво. В източника е посочено, че не се включват елементи със същата стойност.

Например за дърво с корен $10$, ляво поддърво с корен $4$ и дясно поддърво с корен $12$ са валидни само поддървета, които спазват същото правило за подреждане.

\subsection{Търсене}\label{s:18.7.1}

Търсенето на елемент $x$ започва от корена:

\begin{enumerate}
    \item ако дървото е празно, елементът не е намерен;
    \item ако $x$ съвпада със стойността на корена, елементът е намерен;
    \item ако $x$ е по-малък от корена, търсенето продължава в лявото поддърво;
    \item ако $x$ е по-голям от корена, търсенето продължава в дясното поддърво.
\end{enumerate}

\begin{verbatim}
Node* search(Node* tree, int value) {
  if (tree == nullptr || tree->data == value) {
    return tree;
  }

  if (value < tree->data) {
    return search(tree->left, value);
  }

  return search(tree->right, value);
}
\end{verbatim}

Източникът посочва сложност $O(\log n)$ за търсене. Тази оценка се получава при дърво с логаритмична височина. При небалансирано дърво височината може да бъде линейна и тогава търсенето е $O(n)$.

\subsection{Включване}\label{s:18.7.2}

Включването също започва от корена. Ако дървото е празно, се създава нов възел с празни ляво и дясно поддърво. Ако стойността е по-малка от корена, включването продължава в лявото поддърво. Ако е по-голяма, продължава в дясното. При същата стойност включването не се извършва.

\begin{verbatim}
Node* insert(Node* tree, int value) {
  if (tree == nullptr) {
    Node* item = new Node;
    item->data = value;
    item->left = nullptr;
    item->right = nullptr;
    return item;
  }

  if (value < tree->data) {
    tree->left = insert(tree->left, value);
  } else if (value > tree->data) {
    tree->right = insert(tree->right, value);
  }

  return tree;
}
\end{verbatim}

Посочената в източника сложност е $O(\log n)$ при дърво с логаритмична височина. При небалансирано дърво включването може да бъде $O(n)$.

\subsection{Изключване}\label{s:18.7.3}

При изключване първо се търси елементът по правилото за подреденото дърво. Ако елементът не съществува, изключване не се извършва.

Разграничават се следните случаи:

\begin{enumerate}
    \item Изтриваният връх е лист. Той може да бъде премахнат, като съответният указател на неговия баща стане \verb|nullptr|.
    \item Връхът има само ляво поддърво. На негово място се поставя лявото поддърво.
    \item Връхът има само дясно поддърво. На негово място се поставя дясното поддърво.
    \item Връхът има и ляво, и дясно поддърво. Избира се най-левият връх на дясното поддърво. Неговата стойност се поставя на мястото на изтриваната стойност, след което избраният връх се изключва от дясното поддърво.
\end{enumerate}

При последния случай избраният връх е най-малкият елемент в дясното поддърво и следователно запазва свойството за подреждане.

\begin{verbatim}
Node* removeMin(Node* tree, int& value) {
  if (tree->left == nullptr) {
    value = tree->data;
    Node* right = tree->right;
    delete tree;
    return right;
  }

  tree->left = removeMin(tree->left, value);
  return tree;
}

Node* remove(Node* tree, int value) {
  if (tree == nullptr) {
    return nullptr;
  }

  if (value < tree->data) {
    tree->left = remove(tree->left, value);
    return tree;
  }

  if (value > tree->data) {
    tree->right = remove(tree->right, value);
    return tree;
  }

  if (tree->left == nullptr) {
    Node* right = tree->right;
    delete tree;
    return right;
  }

  if (tree->right == nullptr) {
    Node* left = tree->left;
    delete tree;
    return left;
  }

  int successor;
  tree->right = removeMin(tree->right, successor);
  tree->data = successor;
  return tree;
}
\end{verbatim}

Както при търсенето и включването, източникът посочва сложност $O(\log n)$ за изключване. Тя се отнася до дърво с логаритмична височина. В общия случай сложността се определя от височината на дървото.

\subsection{Обхождане на двойно подредено дърво}\label{s:18.7.4}

При смесено обхождане -- ляво поддърво, корен, дясно поддърво -- елементите на двойното подредено дърво се получават във възходящ ред.

При обхождане корен, дясно поддърво, ляво поддърво елементите се получават в низходящ ред. Това следва от правилото, че всички елементи в лявото поддърво са по-малки от корена, а всички в дясното са по-големи.


\subsection{Клас за двойно подредено дърво}\label{s:18.7.5}

Инвариант: за всеки възел $v$ всички стойности в лявото поддърво са строго
по-малки от \verb|v->data|, а всички в дясното --- строго по-големи; равни
стойности не се съхраняват. Смесеното обхождане на класа затова връща
нарастваща редица.

\begin{verbatim}
class SearchTree {
  struct Node {
    int data;
    Node* left;
    Node* right;
  };

  Node* root;

  static void destroy(Node* v) {
    if (v == nullptr) return;
    destroy(v->left);
    destroy(v->right);
    delete v;
  }

  static Node* clone(const Node* v) {
    if (v == nullptr) return nullptr;
    return new Node{v->data, clone(v->left), clone(v->right)};
  }

  static Node* insert(Node* v, int x) {
    if (v == nullptr) return new Node{x, nullptr, nullptr};
    if (x < v->data) v->left = insert(v->left, x);
    else if (x > v->data) v->right = insert(v->right, x);
    return v;             // x == v->data: нищо не се променя
  }

  static Node* minNode(Node* v) {
    while (v->left != nullptr) v = v->left;
    return v;
  }

  static Node* remove(Node* v, int x) {
    if (v == nullptr) return nullptr;
    if (x < v->data) { v->left = remove(v->left, x); return v; }
    if (x > v->data) { v->right = remove(v->right, x); return v; }

    if (v->left == nullptr) {      // нула или един наследник
      Node* r = v->right;
      delete v;
      return r;
    }
    if (v->right == nullptr) {
      Node* l = v->left;
      delete v;
      return l;
    }
    // два наследника: най-малкият вдясно
    Node* s = minNode(v->right);
    v->data = s->data;
    v->right = remove(v->right, s->data);
    return v;
  }

public:
  SearchTree() : root(nullptr) {}
  SearchTree(const SearchTree& other)
      : root(clone(other.root)) {}

  SearchTree& operator=(const SearchTree& other) {
    if (this != &other) {
      Node* copied = clone(other.root);
      destroy(root);
      root = copied;
    }
    return *this;
  }

  ~SearchTree() { destroy(root); }

  bool contains(int x) const {
    const Node* v = root;
    while (v != nullptr) {
      if (x == v->data) return true;
      v = (x < v->data) ? v->left : v->right;
    }
    return false;
  }

  void insert(int x) { root = insert(root, x); }
  void remove(int x) { root = remove(root, x); }
};
\end{verbatim}

Изключването на връх с два наследника заменя стойността му с най-малката в
дясното поддърво и после премахва \emph{нея} --- този възел по построение няма
ляв наследник, затова попада в лесния случай. Заменящата стойност е по-голяма от
всичко вляво и по-малка от всичко останало вдясно, следователно инвариантът се
запазва.

И трите операции --- \verb|contains|, \verb|insert| и \verb|remove| --- слизат
по един-единствен път от корена, затова струват $O(h)$, където $h$ е височината.
Това е $O(\log n)$ само при височина от порядъка на $\log n$; при включване на
вече наредени стойности дървото се изражда в списък и трите операции стават
$O(n)$.

\section{Сравнение на структурите}\label{s:18.8}

Стекът и опашката са линейни структури, но се различават по реда на достъп. При стека се използва LIFO, а при опашката -- FIFO. Списъкът е по-обща линейна структура, защото допуска включване и изключване от произволно място.

Дървото е разклонена структура. При двойното дърво всеки връх има най-много ляв и десен наследник. Физическото представяне чрез указатели следва естествено рекурсивното определение на дървото, докато масивните представяния използват индекси вместо адреси.

\subsection{Сложност на операциите}\label{s:18.8.1}

Таблицата обобщава цената на основните операции. Навсякъде $n$ е броят на
съхранените елементи, а $h$ --- височината на дървото. Оценките са в
\emph{най-лошия} случай, освен изрично отбелязаните като амортизирани. Приема се,
че сравнението и копирането на един елемент струват $O(1)$; при елементи с
по-скъпо копиране цената на преоразмеряването и на \verb|clone| се умножава
съответно.

\begin{center}
\small
\begin{tabular}{llccc}
\toprule
Структура & Реализация & Включване & Изключване & Намиране\\
\midrule
Стек   & статична (масив)      & $O(1)$                    & $O(1)$                    & $O(n)$\\
Стек   & динамична (удвояване) & $O(1)^{1}$                & $O(1)$                    & $O(n)$\\
Стек   & свързана              & $O(1)$                    & $O(1)$                    & $O(n)$\\
\midrule
Опашка & статична (циклична)   & $O(1)$                    & $O(1)$                    & $O(n)$\\
Опашка & динамична             & $O(1)^{1}$                & $O(1)$                    & $O(n)$\\
Опашка & свързана              & $O(1)$                    & $O(1)$                    & $O(n)$\\
\midrule
Списък & едносвързан           & $O(1)^{2}$                & $O(n)^{3}$                & $O(n)$\\
Списък & двусвързан            & $O(1)^{2}$                & $O(1)^{4}$                & $O(n)$\\
Списък & последователен        & $O(n)$                    & $O(n)$                    & $O(n)^{5}$\\
\midrule
Двойно дърво & свързана        & $O(1)^{6}$                & $O(1)^{6}$                & $O(n)$\\
ДПД          & свързана        & $O(h)$                    & $O(h)$                    & $O(h)$\\
\bottomrule
\end{tabular}
\end{center}

\begin{center}
\small
\begin{tabular}{@{}ll@{}}
$^{1}$ амортизирано & $^{2}$ в края; при двусвързан и в началото\\
$^{3}$ по стойност --- търси се предшественикът & $^{4}$ при вече известен възел\\
$^{5}$ $O(\log n)$ при сортиран масив & $^{6}$ при вече известен баща\\
\end{tabular}
\end{center}

Единственият ред, който заблуждава, ако се чете без уговорката, е последният:
$O(h)$ е $\Theta(\log n)$ при балансирано дърво и $\Theta(n)$ при изродено.
Балансирането не се разглежда в този въпрос, затова оценките за ДПД се дават
чрез $h$, а не чрез $\log n$.

\section{Изпитен фокус}\label{s:18.9}

\begin{itemize}
    \item Да се дефинира структура от данни и да се различат логическо описание и физическо представяне.
    \item Да се обяснят стекът и правилото LIFO.
    \item Да се опишат операциите включване, изключване, достъп до върха и търсене в стек.
    \item Да се различат трите реализации, назовани от конспекта: статична (масив с фиксиран капацитет), динамична (масив с удвояване, $O(1)$ амортизирано включване) и свързана (възли с указатели).
    \item Да се сравнят последователното и свързаното представяне на стек.
    \item Да се обяснят опашката и правилото FIFO.
    \item Да се опишат началото и краят на опашката и цикличното последователно представяне.
    \item Да се обясни защо при свързана опашка са нужни указатели към началото и края.
    \item Да се дефинира списък и да се сравнят едносвързаното, двусвързаното и последователното представяне.
    \item Да се покажат връзките в едносвързан и двусвързан списък.
    \item Да се дефинира дърво и двойно дърво рекурсивно.
    \item Да се обяснят корен, наследник, листо, вътрешен връх, ляво и дясно поддърво.
    \item Да се опишат свързаното представяне, верижното представяне чрез масиви и списъкът на бащите.
    \item Да се обясни обхождането и редът на действията при предварително, смесено и последващо обхождане.
    \item Да се дефинира двойно подредено дърво.
    \item Да се възпроизведат алгоритмите за търсене, включване и изключване в двойно подредено дърво.
    \item Да се обясни изключването на връх с нула, един или два наследника.
    \item Да се посочи, че оценките $O(\log n)$ за операциите в двойно подредено дърво предполагат логаритмична височина, докато в общия случай височината може да доведе до линейна сложност.
    \item \textbf{Да се напише клас} за всяка от петте структури --- списък, стек, опашка, двойно кореново дърво и двойно подредено дърво --- с инвариант, конструктор, основни операции и деструктор. Конспектът иска „дефиниране на клас“, а не само свободни функции над отворена структура.
    \item Да се обясни защо копиращ конструктор, операция за присвояване и деструктор вървят заедно при клас, който притежава динамична памет, и какво се случва, ако някой от тях липсва.
    \item Да се възпроизведе таблицата със сложностите (\S18.8.1) заедно с допусканията: най-лош случай, $O(1)$ на елемент за сравнение и копиране, и $O(h)$ вместо $O(\log n)$ за ДПД.
\end{itemize}
